% % % % % % % % % % % % % % % % % % % % % % % % % % % % % % % % % % % % % % % % % % % % % % % % % % % % % % % % % % % % % % % %
% SAT_STYLE_MATH_MCQS.tex
% 1_05.tex
% 
% encoding: UTF-8 
% EOL: LF
%
% format: LaTeX
% indent: spaces (2)
% width: 127
% % % % % % % % % % % % % % % % % % % % % % % % % % % % % % % % % % % % % % % % % % % % % % % % % % % % % % % % % % % % % % % %
\label{1.5}
On considère le produit $P = 3 \times 5 \times 7 \times \cdots \times 97 \times 99$. Parmi ses diviseurs, la plus grande %
puissance de $7$ est:
%
\begin{enumerate}
  \item $10$
  \item $7$
  \item $6$
  \item $11$
  \item $8$
\end{enumerate}
%
\begin{proof}
Il s'agit de trouver le plus grand entier $N$ tel que $7^N \mid P$. 
Remarquons que $7$ est premier: Le lemme d'Euclide nous permet d'affirmer que $7 \mid P = a b$ si et seulement si %
$7 \mid a$ ou $ 7 \mid b$. Attention, cela ne s'applique qu'aux premiers! Pour s'en convaincre, considérons $6$, qui divise %
$24 = (1 \times 2 \times 2) \times 3$ sans diviser ni $1$, ni $2$, ni $3$ (ses diviseurs), ni $2 \times 2 = 4$. 
%
Comme
%
\begin{equation}
P = 3 \times 5 \times 7 \times \cdots \times 97 \times 99, 
\end{equation}
%
il est clair que $7 \mid P$, avec 
%
\begin{equation}
P / 7 = 3 \times 5 \times \mathbf{1} \times 9 \times 11 \times \cdots \times 97 \times 99. 
\end{equation}
%
Peut-on continuer?  Oui, avec $21 = 3 \times 7$. En effet, 
%
\begin{equation}
(P / 7) / 7 = P / 7^2 = 3 \times 5 \times \mathbf{1} \times 9 \times \cdots \times 19 \times 3 \times %
1 \times 23 \times \cdots \times 97 \times 99.
\end{equation}
%
Le prochain impair multiple de $7$ est $35 = 5 \times 7$. D'où:
%
\begin{equation}
P/7^{3} = 3 \times 5 \times \mathbf{1}\times 9 \times \cdots \times 19 \times 3 \times %
1 \times 23 \times \cdots \times 33 \times 5 \times \mathbf{1}\times 37\cdots \times 97 \times 99.
\end{equation}
%
Attention, à l'étape suivante, nous avons $7 \times 7=7^{2}$. Nous divisons \textbf{deux} fois par $7$: 
%
\begin{equation}
P/7^{5}=3 \times 5 \times \mathbf{1}\times 9 \times \cdots \times 33 \times 5 \times %
1 \times 37 \times \cdots \times 47 \times \mathbf{1}\times \mathbf{1}\times \cdots \times 99.
\end{equation}
%
Notre itération s'arrête au moment précis où $99$ est dépassé. Nous prenons donc en compte 
%
\begin{align*} 
63 &= 9 \times 7, \\
77 &= 11 \times 7, \\
91 &= 13 \times 7.
\end{align*}
%
En revanche, $15 \times 7=105>99$ et tous les suivants, notamment $49 \times 7 = 7^3$, sont écartés. En résumé, nous listons %
les multiples impairs de $7$ inférieurs ou égaux à $99$, en notant bien (ci-dessous entre crochets) chaque contribution: 
%
\begin{equation}
7 \{1\}, 21 \{1\}, 35 \{1\}, 49 \{2\}, 63 \{1\}, 77 \{1\}, 91 \{1\}. 
\end{equation}
%
L'entier $N$ recherché est donc le total $1+1+1+2+1+1+1 = 8$. Réponse \textbf{E}.
%
\end{proof}